\section{Introduction}

\xavier{J'ai passé en revue les traductions, indépendamment du message envoyé le 18/05 concernant le fond de l'introduction}

The ability to generate several variants of the same object automatically is
becoming increasingly frequent in many application areas. %
The CAO field is a typical example which uses history-based parametric
systems. %
They are dedicated to the recording of construction processes, which are often
both tedious and time-consuming. Then, they allow to automatically re-evaluate
those processes by modifying some of their parameters, in order to generate
different versions of the object to produce. %
By the same token, in the field of archaeology, remaining data found on the
working field often represent only vestiges of ancients buildings. %
Therefore, by means of 3D reproduction, archaeologists develop a number of
hypotheses they expect to test while being able to quickly model and visualise
them \cite{Quattrini_Baleani_2015}. %
Grammar-based procedural methods are commonly used to generate variations of a
building \cite{haegler2009procedural}. %
However, those grammars are based on specific corpus of information which are
difficult to transpose to other case studies. %
In this paper, we propose to use a rule-based graph transformation formalism,
and more specifically the Jerboa \cite{belhaouari2014jerboa} library, to make
the development of dedicated modellers for specific applications easier. %
Rule-based formalisms are commonly used in a number of fields such as the
modelling of plant growth with L-systems \cite{lindenmayer1974adding}, of wood's
internal structure \cite{terraz20093gmap}, of architectural buildings
\cite{muller2006procedural} or of virtual cities 3D models
\cite{esri2023cityengine, beirao2010implementing}. %
Contrary to other approaches, Jerboa is independent from any specific
application field and does not require ad hoc operations to be manually coded. %
Operations are formalised by the construction of rules within the Jerboa
interface, thus enabling their quick development. It also ensures the
topological constistency of the underlying geometric model regardless of the
applied operations []. %
Jerboa is based on the "Generalized Maps" (or "G-maps") topological model
\cite{lienhardt1991topological, Damiand-Lienhardt14}. This model matches a
specific class of labelled graph and allows the homogeneous modelling of
quasi-manifolds in any dimension. %
Number of applications already make use of Jerboa and/or G-maps in fields such
as architecture [], geology \cite{gauthier2016topological}
(\cite{halbwachs2000generalized}?) or physics-based modelling
\cite{salah2017general}. %
Although Jerboa provides a number of advantages, it does not provide the
mechanisms required to quickly generate variations from a base model. As a
matter of fact, CAO's history-based parametric systems provide such
mechanisms. %
More precisely, parametric systems can be thought as a dual structure with, on
the one hand, the geometric model corresponding to the modelled object and, on
the other hand, the successive operations (and their parameters) of the
modelling process. %
The topological parameters of an operation reference topological entities
(vertices, edges, faces, and so on) defined prior to the modelling process (for
example, in Fig.~\ref{fig:persistent-naming-pb-eval}, a groove is applied onto
previously defined face $f_{1}$). %
The parametric specification can be slightly modified and re-evaluated to
generate a variant of the resulting object. Still, such an approach requires
from the system to guarantee the \textit{persistence} of the topological
references. %
%
\begin{figure}[ht]
  \centering
  \begin{subfigure}[t]{\linewidth}
    \centering
    \includegraphics[width=.8\textwidth]{persistent-naming-pb-1}
    \caption{Initial evaluation}
    \label{fig:persistent-naming-pb-eval}
  \end{subfigure}
  \begin{subfigure}[t]{\linewidth}
    \centering
    \includegraphics[width=.8\textwidth]{persistent-naming-pb-2}
    \caption{Re-Evaluation}
    \label{fig:persistent-naming-pb-reval}
  \end{subfigure}
  \caption{Parametric specification}
  \label{fig:persistent-naming-pb}
\end{figure}
%
This issue, known as \textit{persistent naming}, is illustrated in
Fig.~\ref{fig:persistent-naming-pb}. %
%
The initial parametric specification consists of three constructive operations
(Fig.~\ref{fig:persistent-naming-pb-eval}): (1) CreateCylinder($geo\_param1$);
(2) RoundedGroove($f_{1}$, $geo\_param2$); (3) CylindricalProtrusion($f_{2}$,
$geo\_param3$). %
%
When identifiers or pointers to the topological parameters of a geometric model
(\textit{e.g.} the identifier of the face $f_{2}$ as a parameter of the
CylindricalProtrusion operation) are used as topological parameters, the issue
of the persistence of these references at re-evaluation comes up. %
%
For example, in Fig.~\ref{fig:persistent-naming-pb-reval}, the rounded groove's
length is reduced. %
%
The face $f_{1}$ is not split anymore, unlike during the initial evaluation.
Thus, faces $f_{2}$ and $f_{3}$ are not created anymore, the identifiers and
pointers to the entities are obviously different and, therefore, the cylindrical
protrusion can no longer be re-applied onto $f_{2}$. %
%
Hence the necessity to use persistent identifiers as operations' parameters,
which make possible to unambiguously characterise entities and find their match
at re-evaluation. In Fig.~\ref{fig:persistent-naming-pb}, using a persistent name
to characterise the face $f_{2}$ during the initial evaluation allows to match
it with the face $f_{x}$ at re-evaluation. %
%
The persistent naming problem has been studied for decades in the CAO's field
\cite{marcheix2002survey,farjana2018mechanisms} (\cite{krispel2015survey}?), yet
none has been written jointly with the use of graph transformation rules.
\xavier{Il faut quand même justifier à quoi les règles de transformation peuvent
  servir : quel est le problème qu'elles pourraient résoudre, et qui n'est pas
  géré par les autres approches de la littérature ?} %
We propose: Firstly, to widen the naming problem studies to the rule-based graph
transformation modelling systems; Secondly, to integrate the mechanisms of
re-evaluation for parametric system within Jerboa. %
To this end, we propose a persistent naming method taking advantage of the
rules' formalisation of operations and their ability to precisely describe the
history of topological entities such that these entities are uniquely and
homogeneously characterised for all dimensions. %
Most existing methods require the tracking of numerous topological entities and
consider the persistent naming problem only from the parameters' modifications
of a parametric specification standpoint\cite{kripac1995mechanism,
  chen1995editability, wu2001face}. %
Our solution tracks only the entities used in the parametric specification and
the ones they originate from. %
Moreover, not only the naming problem is tackled within the usual framework of
specification's parameters edition but we also take the specification edition
(\textit{i.e.} adding, deleting and changing the order of operations inside the
specification) into account. %
In section 2, we present the main concepts necessary to carry out the persistent
naming mechanisms within the framework of a system, making use of graph
transformation rules. %
We focus on both G-maps and Jerboa's rules. We also describe how to
automatically detect topological changes (creation, deletion, split, merging,
modification) that may occur onto topological entities upon a rule's
application. %
%
Section 3 is dedicated to structures defined for parametric specification and
persistent naming. %
Then, Sections 4 and 5 tackle the concepts of evaluation and re-evaluation
Direct Acyclic Graphs (or DAGs) enabling the tracking of the topological
entities' evolutions, the entities they originate from and their matching with
the corresponding entities at re-evaluation. %
